% Options for packages loaded elsewhere
% Options for packages loaded elsewhere
\PassOptionsToPackage{unicode}{hyperref}
\PassOptionsToPackage{hyphens}{url}
\PassOptionsToPackage{dvipsnames,svgnames,x11names}{xcolor}
%
\documentclass[
  spanish,
  12pt,
  a4paper,
]{article}
\usepackage{xcolor}
\usepackage[margin=2.5cm]{geometry}
\usepackage{amsmath,amssymb}
\setcounter{secnumdepth}{-\maxdimen} % remove section numbering
\usepackage{iftex}
\ifPDFTeX
  \usepackage[T1]{fontenc}
  \usepackage[utf8]{inputenc}
  \usepackage{textcomp} % provide euro and other symbols
\else % if luatex or xetex
  \usepackage{unicode-math} % this also loads fontspec
  \defaultfontfeatures{Scale=MatchLowercase}
  \defaultfontfeatures[\rmfamily]{Ligatures=TeX,Scale=1}
\fi
\usepackage{lmodern}
\ifPDFTeX\else
  % xetex/luatex font selection
  \setmainfont[]{Times New Roman}
\fi
% Use upquote if available, for straight quotes in verbatim environments
\IfFileExists{upquote.sty}{\usepackage{upquote}}{}
\IfFileExists{microtype.sty}{% use microtype if available
  \usepackage[]{microtype}
  \UseMicrotypeSet[protrusion]{basicmath} % disable protrusion for tt fonts
}{}
\usepackage{setspace}
\makeatletter
\@ifundefined{KOMAClassName}{% if non-KOMA class
  \IfFileExists{parskip.sty}{%
    \usepackage{parskip}
  }{% else
    \setlength{\parindent}{0pt}
    \setlength{\parskip}{6pt plus 2pt minus 1pt}}
}{% if KOMA class
  \KOMAoptions{parskip=half}}
\makeatother
% Make \paragraph and \subparagraph free-standing
\makeatletter
\ifx\paragraph\undefined\else
  \let\oldparagraph\paragraph
  \renewcommand{\paragraph}{
    \@ifstar
      \xxxParagraphStar
      \xxxParagraphNoStar
  }
  \newcommand{\xxxParagraphStar}[1]{\oldparagraph*{#1}\mbox{}}
  \newcommand{\xxxParagraphNoStar}[1]{\oldparagraph{#1}\mbox{}}
\fi
\ifx\subparagraph\undefined\else
  \let\oldsubparagraph\subparagraph
  \renewcommand{\subparagraph}{
    \@ifstar
      \xxxSubParagraphStar
      \xxxSubParagraphNoStar
  }
  \newcommand{\xxxSubParagraphStar}[1]{\oldsubparagraph*{#1}\mbox{}}
  \newcommand{\xxxSubParagraphNoStar}[1]{\oldsubparagraph{#1}\mbox{}}
\fi
\makeatother


\usepackage{longtable,booktabs,array}
\usepackage{calc} % for calculating minipage widths
% Correct order of tables after \paragraph or \subparagraph
\usepackage{etoolbox}
\makeatletter
\patchcmd\longtable{\par}{\if@noskipsec\mbox{}\fi\par}{}{}
\makeatother
% Allow footnotes in longtable head/foot
\IfFileExists{footnotehyper.sty}{\usepackage{footnotehyper}}{\usepackage{footnote}}
\makesavenoteenv{longtable}
\usepackage{graphicx}
\makeatletter
\newsavebox\pandoc@box
\newcommand*\pandocbounded[1]{% scales image to fit in text height/width
  \sbox\pandoc@box{#1}%
  \Gscale@div\@tempa{\textheight}{\dimexpr\ht\pandoc@box+\dp\pandoc@box\relax}%
  \Gscale@div\@tempb{\linewidth}{\wd\pandoc@box}%
  \ifdim\@tempb\p@<\@tempa\p@\let\@tempa\@tempb\fi% select the smaller of both
  \ifdim\@tempa\p@<\p@\scalebox{\@tempa}{\usebox\pandoc@box}%
  \else\usebox{\pandoc@box}%
  \fi%
}
% Set default figure placement to htbp
\def\fps@figure{htbp}
\makeatother



\ifLuaTeX
\usepackage[bidi=basic]{babel}
\else
\usepackage[bidi=default]{babel}
\fi
\ifPDFTeX
\else
\babelfont{rm}[]{Times New Roman}
\fi
% get rid of language-specific shorthands (see #6817):
\let\LanguageShortHands\languageshorthands
\def\languageshorthands#1{}


\setlength{\emergencystretch}{3em} % prevent overfull lines

\providecommand{\tightlist}{%
  \setlength{\itemsep}{0pt}\setlength{\parskip}{0pt}}



 


\usepackage[utf8]{inputenc}
\usepackage{textgreek}
\DeclareUnicodeCharacter{00B5}{\textmu}
\usepackage{lineno}
\linenumbers
\usepackage{booktabs}
\usepackage{longtable}
\usepackage{setspace}
\onehalfspacing
\usepackage{fancyhdr}
\pagestyle{fancy}
\fancyhf{}
\fancyhead[L]{TFG · v0.1 · 2026-04-20}
\fancyhead[R]{\thepage}
\fancyfoot[C]{\small Borrador preliminar — no distribuir}
\usepackage{csquotes}
\usepackage{enumitem}
\setlist{nosep, leftmargin=1.5em}
\makeatletter
\@ifpackageloaded{caption}{}{\usepackage{caption}}
\AtBeginDocument{%
\ifdefined\contentsname
  \renewcommand*\contentsname{Tabla de contenidos}
\else
  \newcommand\contentsname{Tabla de contenidos}
\fi
\ifdefined\listfigurename
  \renewcommand*\listfigurename{Listado de Figuras}
\else
  \newcommand\listfigurename{Listado de Figuras}
\fi
\ifdefined\listtablename
  \renewcommand*\listtablename{Listado de Tablas}
\else
  \newcommand\listtablename{Listado de Tablas}
\fi
\ifdefined\figurename
  \renewcommand*\figurename{Figura}
\else
  \newcommand\figurename{Figura}
\fi
\ifdefined\tablename
  \renewcommand*\tablename{Tabla}
\else
  \newcommand\tablename{Tabla}
\fi
}
\@ifpackageloaded{float}{}{\usepackage{float}}
\floatstyle{ruled}
\@ifundefined{c@chapter}{\newfloat{codelisting}{h}{lop}}{\newfloat{codelisting}{h}{lop}[chapter]}
\floatname{codelisting}{Listado}
\newcommand*\listoflistings{\listof{codelisting}{Listado de Listados}}
\makeatother
\makeatletter
\makeatother
\makeatletter
\@ifpackageloaded{caption}{}{\usepackage{caption}}
\@ifpackageloaded{subcaption}{}{\usepackage{subcaption}}
\makeatother
\usepackage{bookmark}
\IfFileExists{xurl.sty}{\usepackage{xurl}}{} % add URL line breaks if available
\urlstyle{same}
\hypersetup{
  pdftitle={Cinética de seronegativización por IgG-ELISA tras ivermectina en estrongiloidiasis: cohorte retrospectiva (HCDGU, 2015--2025)},
  pdfauthor={Yan Rodríguez Hachimaru},
  pdflang={es},
  colorlinks=true,
  linkcolor={blue},
  filecolor={Maroon},
  citecolor={Blue},
  urlcolor={Blue},
  pdfcreator={LaTeX via pandoc}}


\title{Cinética de seronegativización por IgG-ELISA tras ivermectina en
estrongiloidiasis: cohorte retrospectiva (HCDGU, 2015--2025)}
\usepackage{etoolbox}
\makeatletter
\providecommand{\subtitle}[1]{% add subtitle to \maketitle
  \apptocmd{\@title}{\par {\large #1 \par}}{}{}
}
\makeatother
\subtitle{Trabajo Fin de Grado · Facultad de Medicina y Ciencias de la
Salud, UAH}
\author{Yan Rodríguez Hachimaru}
\date{2026-04-20}
\begin{document}
\maketitle


\setstretch{1.5}
\section*{Cinética de seronegativización por IgG-ELISA tras ivermectina
en estrongiloidiasis: cohorte retrospectiva (HCDGU,
2015--2025)}\label{cinuxe9tica-de-seronegativizaciuxf3n-por-igg-elisa-tras-ivermectina-en-estrongiloidiasis-cohorte-retrospectiva-hcdgu-20152025}
\addcontentsline{toc}{section}{Cinética de seronegativización por
IgG-ELISA tras ivermectina en estrongiloidiasis: cohorte retrospectiva
(HCDGU, 2015--2025)}

\textbf{Autor:} Yan Rodríguez Hachimaru. \textbf{Tutor/a:}
{[}PENDIENTE{]}. \textbf{Filiación:} Facultad de Medicina y Ciencias de
la Salud, Universidad de Alcalá, Alcalá de Henares (Madrid), España.
\textbf{Centro del estudio:} Hospital Central de la Defensa ``Gómez
Ulla'' (HCDGU), Madrid, España. \textbf{Correspondencia:} {[}email
PENDIENTE{]}.

\textbf{Palabras clave (ES):} estrongiloidiasis; \emph{Strongyloides
stercoralis}; ivermectina; seronegativización; IgG-ELISA; Kaplan--Meier;
modelo de Cox; cinética serológica; cohorte retrospectiva.

\begin{center}\rule{0.5\linewidth}{0.5pt}\end{center}

\section*{Serological clearance kinetics of IgG-ELISA after ivermectin
in strongyloidiasis: a retrospective cohort (HCDGU,
2015--2025)}\label{serological-clearance-kinetics-of-igg-elisa-after-ivermectin-in-strongyloidiasis-a-retrospective-cohort-hcdgu-20152025}
\addcontentsline{toc}{section}{Serological clearance kinetics of
IgG-ELISA after ivermectin in strongyloidiasis: a retrospective cohort
(HCDGU, 2015--2025)}

\textbf{Author:} Yan Rodríguez Hachimaru. \textbf{Supervisor:}
{[}PENDING{]}. \textbf{Affiliation:} School of Medicine and Health
Sciences, Universidad de Alcalá, Alcalá de Henares (Madrid), Spain.
\textbf{Study site:} Hospital Central de la Defensa ``Gómez Ulla''
(HCDGU), Madrid, Spain. \textbf{Correspondence:} {[}email PENDING{]}.

\textbf{Keywords (EN):} strongyloidiasis; \emph{Strongyloides
stercoralis}; ivermectin; seroreversion; IgG-ELISA; Kaplan--Meier; Cox
regression; serological kinetics; retrospective cohort.

\newpage

\subsection*{Resumen}\label{resumen}
\addcontentsline{toc}{subsection}{Resumen}

\textbf{Introducción.} La estrongiloidiasis humana, causada por
\emph{Strongyloides stercoralis}, es una helmintiasis crónica con
potencial letal en inmunodeprimidos cuya monitorización pos-tratamiento
se realiza mediante serología IgG-ELISA. La cinética del descenso
serológico tras ivermectina en cohortes no endémicas no está plenamente
caracterizada.

\textbf{Objetivos.} Describir el tiempo hasta la seronegativización
(IgG-ELISA \textless{} 1,1) en adultos tratados con ivermectina e
identificar predictores basales asociados.

\textbf{Material y métodos.} Estudio observacional retrospectivo
(2015--2025) en el Hospital Central de la Defensa ``Gómez Ulla''
(Madrid). Se incluyeron 44 adultos con estrongiloidiasis confirmada y
tratados con ivermectina 200 μg/kg. Se construyeron curvas de
Kaplan--Meier, modelos de Cox univariables (ajuste Benjamini--Hochberg
False Discovery Rate, BH--FDR) y un modelo lineal mixto (LMM) sobre
log(IgG).

\textbf{Resultados.} De 44 pacientes, 39 aportaron datos evaluables. La
mediana de tiempo hasta seronegativización fue 8,8 meses (IC 95 \%
4,3--11,4), con 36 \% a 6 meses y 88 \% a 24 meses. En Cox univariable,
la IgG basal ≥ mediana se asoció a menor HR de aclaramiento (HR 0,19;
p=0,004) y el sexo masculino a mayor HR (HR 3,29; p=0,007); ningún
predictor sobrevivió al ajuste BH--FDR. El LMM estimó pendiente −0,068
unidades/mes (p\textless0,001), vida media serológica 10,2 meses.

\textbf{Conclusiones.} La seronegativización es tardía (mediana
\textasciitilde9 meses) y predecible por cinética de primer orden. La
IgG basal elevada es un marcador exploratorio de aclaramiento más lento
que requiere validación en cohortes mayores.

\begin{center}\rule{0.5\linewidth}{0.5pt}\end{center}

\subsection*{Abstract}\label{abstract}
\addcontentsline{toc}{subsection}{Abstract}

\textbf{Background.} Human strongyloidiasis, caused by
\emph{Strongyloides stercoralis}, is a chronic helminth infection with
potentially fatal outcomes in immunocompromised patients and is
followed-up serologically by IgG-ELISA. Serological clearance kinetics
after ivermectin in non-endemic cohorts remain incompletely
characterised.

\textbf{Objectives.} To describe time to seroreversion (IgG-ELISA
\textless{} 1.1) in adults treated with ivermectin and to identify
associated baseline predictors.

\textbf{Methods.} Retrospective observational study (2015--2025) at
Hospital Central de la Defensa ``Gómez Ulla'' (Madrid, Spain). We
included 44 adults with confirmed strongyloidiasis treated with
ivermectin 200 μg/kg. We used Kaplan--Meier estimates, univariate Cox
models (Benjamini--Hochberg False Discovery Rate adjustment, BH--FDR),
and a linear mixed-effects model (LMM) on log(IgG).

\textbf{Results.} Of 44 patients, 39 contributed evaluable follow-up.
Median time to seroreversion was 8.8 months (95 \% CI 4.3--11.4); 36 \%
at 6 months and 88 \% at 24 months. In univariate Cox, baseline IgG ≥
median was associated with slower clearance (HR 0.19; p=0.004) and male
sex with faster clearance (HR 3.29; p=0.007); no predictor remained
significant after BH--FDR adjustment. LMM estimated a slope of −0.068
units/month (p\textless0.001) with a serological half-life of 10.2
months.

\textbf{Conclusions.} Seroreversion is delayed (median \textasciitilde9
months) and follows first-order kinetics. High baseline IgG is an
exploratory marker of slower clearance warranting replication.

\newpage

\subsection*{Abreviaturas y glosario}\label{abreviaturas-y-glosario}
\addcontentsline{toc}{subsection}{Abreviaturas y glosario}

\begin{longtable}[]{@{}
  >{\raggedright\arraybackslash}p{(\linewidth - 2\tabcolsep) * \real{0.5000}}
  >{\raggedright\arraybackslash}p{(\linewidth - 2\tabcolsep) * \real{0.5000}}@{}}
\toprule\noalign{}
\begin{minipage}[b]{\linewidth}\raggedright
Abreviatura
\end{minipage} & \begin{minipage}[b]{\linewidth}\raggedright
Definición
\end{minipage} \\
\midrule\noalign{}
\endhead
\bottomrule\noalign{}
\endlastfoot
BH--FDR & Benjamini--Hochberg False Discovery Rate (ajuste por
comparaciones múltiples) \\
CNM & Centro Nacional de Microbiología (Instituto de Salud Carlos
III) \\
CSL & Citation Style Language \\
ELISA & Enzyme-Linked ImmunoSorbent Assay \\
HCDGU & Hospital Central de la Defensa ``Gómez Ulla'' \\
HR & Hazard Ratio \\
IC & Intervalo de Confianza \\
IgG & Inmunoglobulina G \\
IQR & Rango intercuartílico (\emph{Interquartile Range}) \\
IVM & Ivermectina \\
KM & Kaplan--Meier \\
LMM & Linear Mixed Model (modelo lineal mixto) \\
MICE & \emph{Multiple Imputation by Chained Equations} \\
MyM & Material y Métodos \\
PCR & Reacción en cadena de la polimerasa \\
STROBE & \emph{Strengthening the Reporting of Observational Studies in
Epidemiology} \\
TFG & Trabajo Fin de Grado \\
UAH & Universidad de Alcalá \\
\end{longtable}

\newpage

\subsection*{Introducción}\label{introducciuxf3n}
\addcontentsline{toc}{subsection}{Introducción}

\subsubsection{Contexto
epidemiológico}\label{contexto-epidemioluxf3gico}

La estrongiloidiasis humana es una infección crónica causada por el
nematodo \emph{Strongyloides stercoralis}, con una prevalencia global
estimada de 300--600 millones de personas, principalmente en regiones
tropicales y subtropicales. {[}CITE?:buonfrate\_2023\_cmr{]}

En Europa y en otros países no endémicos, la enfermedad se diagnostica
predominantemente en población migrante procedente de América Latina,
África subsahariana y Asia. {[}CITED:asundi\_2019\_lancetgh{]}

Series españolas multicéntricas recientes describen una cohorte
predominantemente masculina (54--72 \%) y de origen sudamericano, con
una prevalencia de inmunosupresión basal del 6,4 \%.
{[}CITED:salvador\_2019\_redivi{]}

La presentación clínica más frecuente es asintomática o inespecífica, lo
que obliga a estrategias de cribado serológico en poblaciones de riesgo.
{[}CITED:salvador\_2020\_systematic{]}

\subsubsection{Justificación clínica}\label{justificaciuxf3n-cluxednica}

\emph{S. stercoralis} puede persistir décadas en el hospedador mediante
auto-infección endógena, con riesgo de hiperinfestación y
estrongiloidiasis diseminada en pacientes inmunodeprimidos (letalidad
\textgreater{} 60 \%). {[}CITED:czeresnia\_2022\_lung{]}

La ivermectina a 200 μg/kg en dosis única es el tratamiento de elección
en el paciente inmunocompetente, con tasas de curación \textgreater{} 90
\% en ensayos controlados. {[}CITED:buonfrate\_2019\_strongtreat{]}

La confirmación de curación tras ivermectina se basa habitualmente en
serología IgG-ELISA, dado que el rendimiento parasitológico
(coprocultivo, Baermann, PCR) es bajo en infecciones crónicas de baja
carga. {[}CITED:buonfrate\_2023\_cmr{]}

\subsubsection{Problema a abordar}\label{problema-a-abordar}

La cinética del descenso serológico post-tratamiento no está
completamente caracterizada en cohortes europeas no endémicas.
{[}CITE?:repetto\_2018\_longterm{]}

Los estudios publicados aportan tiempos medianos de seronegativización
heterogéneos (6--24 meses) y emplean umbrales de ``negatividad''
variables según el fabricante del kit y la definición operativa del
equipo asistencial. {[}CITE?:bisoffi\_2014\_salvador{]}

Los predictores basales del tiempo hasta seronegativización (edad, sexo,
título basal, eosinofilia, inmunosupresión, carga parasitológica) han
sido evaluados de forma dispersa y sin correción por comparaciones
múltiples en cohortes pequeñas. {[}CITE?:gamez\_2022\_spanish{]}

Existe por tanto una necesidad clínica de describir con precisión la
dinámica de aclaramiento serológico en la práctica real española y de
identificar, en modo exploratorio, marcadores basales que informen el
protocolo de seguimiento. {[}NO\_CITE{]}

\subsubsection{Objetivos}\label{objetivos}

\textbf{Objetivo principal.} Caracterizar el tiempo hasta la
seronegativización absoluta (IgG-ELISA \textless{} 1,1) en una cohorte
retrospectiva de adultos tratados con ivermectina en el Hospital Central
de la Defensa ``Gómez Ulla'' (Madrid, 2015--2025). {[}NO\_CITE{]}

\textbf{Objetivos secundarios.}

\begin{enumerate}
\def\labelenumi{\arabic{enumi}.}
\tightlist
\item
  Estimar tasas de seronegativización a 6, 12, 18 y 24 meses mediante
  análisis de Kaplan--Meier. {[}NO\_CITE{]}
\item
  Explorar predictores basales (demográficos, inmunológicos,
  serológicos, parasitológicos, terapéuticos) mediante modelos de Cox
  univariables con ajuste Benjamini--Hochberg False Discovery Rate
  (BH--FDR). {[}NO\_CITE{]}
\item
  Cuantificar la pendiente de descenso de log(IgG) y la vida media
  serológica mediante un modelo lineal mixto (LMM) con cinética de
  primer orden. {[}NO\_CITE{]}
\item
  Describir la persistencia serológica y los casos de recidiva sin
  evidencia parasitológica. {[}NO\_CITE{]}
\end{enumerate}

\textbf{Encuadre.} Estudio piloto exploratorio, no diseñado para
inferencia causal; los hallazgos se interpretan como generadores de
hipótesis para cohortes mayores. {[}NO\_CITE{]}

\subsection*{Material y métodos}\label{material-y-muxe9todos}
\addcontentsline{toc}{subsection}{Material y métodos}

\subsubsection{Diseño y ámbito}\label{diseuxf1o-y-uxe1mbito}

Se realizó un estudio observacional retrospectivo en el Hospital Central
de la Defensa ``Gómez Ulla'' (HCDGU, Madrid, España) sobre el periodo
2015--2025. {[}NO\_CITE{]}

Se incluyeron adultos (≥18 años) con diagnóstico confirmado de \emph{S.
stercoralis} y tratamiento con ivermectina oral 200 μg/kg (dosis única o
pauta múltiple en casos de inmunosupresión o carga parasitaria elevada).
{[}NO\_CITE{]}

El seguimiento protocolizado contemplaba determinaciones basales y
controles post-tratamiento a los 3, 6, 12, 18 y 24 meses. {[}NO\_CITE{]}

El estudio siguió las recomendaciones STROBE para estudios
observacionales. {[}CITED:vonelm\_2007\_strobe{]}

\subsubsection{Diagnóstico y definición
operativa}\label{diagnuxf3stico-y-definiciuxf3n-operativa}

El diagnóstico se estableció por al menos uno de los siguientes
criterios: (a) coprocultivo positivo, (b) PCR positiva en heces, o (c)
serología IgG-ELISA \textgreater{} 1,1 con clínica o epidemiología
compatible. {[}NO\_CITE{]}

La serología se realizó mediante IgG-ELISA comercial
(\emph{Strongyloides ratti} antigen-based) en el CNM (Centro Nacional de
Microbiología) y en el servicio de Microbiología del HCDGU.
{[}CITE?:bisoffi\_2014\_salvador{]}

Se registró el cambio de kit ELISA durante el periodo de estudio como
posible fuente de variabilidad técnica. {[}NO\_CITE{]}

La ausencia de IgG basal en el 18 \% de la muestra (8 de 44) se
consideró aleatoria (prueba de Little, p = 0,11); no se aplicaron
técnicas de imputación múltiple. {[}NO\_CITE{]}

\subsubsection{Exposición y endpoints}\label{exposiciuxf3n-y-endpoints}

La exposición principal fue la administración de ivermectina 200 μg/kg
(mono-dosis o pautas múltiples). {[}NO\_CITE{]}

El \textbf{criterio de valoración primario} fue la seronegativización
absoluta, definida como el primer descenso de IgG-ELISA por debajo de
1,1 en cualquier determinación post-tratamiento. {[}NO\_CITE{]}

Los criterios de valoración secundarios (análisis de sensibilidad)
incluyeron:

\begin{itemize}
\tightlist
\item
  Seronegativización con umbral más estricto (\textless{} 0,9).
  {[}NO\_CITE{]}
\item
  Seronegativización confirmada en dos determinaciones consecutivas
  (×2). {[}NO\_CITE{]}
\item
  Seronegativización relativa (descenso ≥ 40 \% desde basal), según el
  umbral Salvador. {[}CITE?:bisoffi\_2014\_salvador{]}
\item
  Normalización eosinofílica (\textless{} 0,5 × 10³/μL). {[}NO\_CITE{]}
\item
  Seronegativización parasitológica (dos coprocultivos o PCR negativos
  consecutivos). {[}NO\_CITE{]}
\end{itemize}

\subsubsection{Covariables basales}\label{covariables-basales}

Se registraron: edad al diagnóstico, sexo, origen geográfico
(clasificado por región), inmunodepresión basal, comorbilidad
significativa, IgG basal (continua y categorizada por
mediana/P75/terciles), eosinofilia basal (continua y umbral ≥ 0,5 ×
10³/μL), resultado parasitológico basal, pauta terapéutica (mono vs
múltiple) y retratamiento. {[}NO\_CITE{]}

\subsubsection{Análisis estadístico}\label{anuxe1lisis-estaduxedstico}

Las variables continuas se expresaron como mediana (rango
intercuartílico, IQR) y las categóricas como frecuencias absolutas y
relativas. {[}NO\_CITE{]}

Se construyeron curvas de Kaplan--Meier (KM) con intervalo de confianza
95 \% calculado por el método de Greenwood.
{[}CITE?:kaplan\_meier\_1958{]}

Los predictores basales se evaluaron mediante modelos de Cox
univariables sobre 15 covariables, con ajuste por comparaciones
múltiples mediante el procedimiento Benjamini--Hochberg False Discovery
Rate (BH--FDR, α = 0,05). {[}CITED:benjamini\_1995\_fdr{]}

Un modelo de Cox multivariable estratificado por sexo se ajustó con las
variables que presentaron p \textless{} 0,05 en el análisis univariable,
evaluando el estadístico C de Harrell. {[}CITE?:harrell\_1996\_c{]}

La cinética paramétrica de descenso serológico se modeló mediante un LMM
sobre log(IgG) con intercepto aleatorio por paciente y tiempo como
efecto fijo, asumiendo cinética de primer orden.
{[}CITE?:laird\_ware\_1982\_lmm{]}

Se exploró la interacción tiempo × sexo en el LMM. {[}NO\_CITE{]}

La vida media serológica se derivó como t₁/₂ = ln(2) / pendiente.
{[}NO\_CITE{]}

Se consideró significación estadística para p \textless{} 0,05
(bilateral). {[}NO\_CITE{]}

El análisis se realizó en Python 3.14 con las librerías \texttt{pandas},
\texttt{statsmodels} y \texttt{lifelines}; el código completo es
reproducible a partir del repositorio del proyecto. {[}NO\_CITE{]}

\subsubsection{Aspectos éticos}\label{aspectos-uxe9ticos}

El protocolo fue evaluado por el Comité de Ética de la Investigación con
Medicamentos (CEIM) del HCDGU (referencia PENDIENTE). {[}NO\_CITE{]}

Al tratarse de un estudio retrospectivo con datos anonimizados, se
solicitó la exención del consentimiento informado individual, de acuerdo
con el Reglamento (UE) 2016/679 y la Ley Orgánica 3/2018. {[}NO\_CITE{]}

\subsection*{Resultados}\label{resultados}
\addcontentsline{toc}{subsection}{Resultados}

\subsubsection{Flujo de pacientes y descripción de la
cohorte}\label{flujo-de-pacientes-y-descripciuxf3n-de-la-cohorte}

Del total de 198 registros cribados, 84 cumplieron criterios
diagnósticos confirmados y 64 recibieron ivermectina; tras descartar
pacientes sin seguimiento o sin diagnóstico en HCDGU, la cohorte final
incluyó 44 pacientes (retención 22,2 \%) (Figura 1). {[}NO\_CITE{]}

La cohorte fue mayoritariamente de origen sudamericano (81,8 \%) y
diagnosticada mediante serología (86,4 \%); 52,3 \% recibió ivermectina
en dosis única. {[}NO\_CITE{]}

La mediana de seguimiento serológico fue 11,7 meses (IQR 3,1--22,2).
{[}NO\_CITE{]}

Las características basales detalladas se presentan en la \textbf{Tabla
1} (ver \texttt{manuscrito/figuras/tab-basales.md}). {[}NO\_CITE{]}

\subsubsection{Cinética de
seronegativización}\label{cinuxe9tica-de-seronegativizaciuxf3n}

En los 39 pacientes evaluables para análisis de supervivencia
serológica, la mediana de tiempo hasta seronegativización absoluta
(IgG-ELISA \textless{} 1,1) fue \textbf{8,8 meses} (IC 95 \% 4,3--11,4).
{[}NO\_CITE{]}

Las proporciones acumuladas de seronegativización fueron 36 \% a los 6
meses, 61 \% a los 12 meses, 78 \% a los 18 meses y 88 \% a los 24 meses
(\textbf{Figura 2}). {[}NO\_CITE{]}

El análisis de sensibilidad con umbral más estricto (\textless{} 0,9)
desplazó la mediana a \textbf{10,8 meses} (reducción de eventos de 27 a
25), confirmando la robustez del hallazgo. {[}NO\_CITE{]}

La seronegativización confirmada en dos determinaciones consecutivas
presentó una mediana de \textbf{16,4 meses}. {[}NO\_CITE{]}

La normalización eosinofílica fue significativamente más precoz (mediana
\textbf{3,0 meses}) (\textbf{Figura 3}). {[}NO\_CITE{]}

\subsubsection{Predictores basales (Cox
univariable)}\label{predictores-basales-cox-univariable}

En el análisis de Cox univariable sobre 15 covariables, dos variables
mostraron las asociaciones más fuertes con el tiempo de
seronegativización (\textbf{Figura 4A}, \textbf{Tabla 2}):

\begin{itemize}
\tightlist
\item
  IgG basal ≥ mediana: HR 0,19 (IC 95 \% 0,06--0,59; p = 0,004).
  {[}NO\_CITE{]}
\item
  Sexo masculino: HR 3,29 (IC 95 \% 1,38--7,86; p = 0,007).
  {[}NO\_CITE{]}
\end{itemize}

Tras el ajuste Benjamini--Hochberg False Discovery Rate (BH--FDR),
ninguna covariable mantuvo significación (p-ajustada \textgreater{}
0,05). {[}NO\_CITE{]}

Las curvas de Kaplan--Meier estratificadas por estas variables
(\textbf{Figura 4B--D}) visualizan la separación de trayectorias.
{[}NO\_CITE{]}

\subsubsection{Modelo multivariable}\label{modelo-multivariable}

El modelo de Cox multivariable estratificado por sexo (IgG basal elevada
+ eosinofilia basal como covariables) alcanzó un estadístico C de
Harrell de \textbf{0,676} (\textbf{Tabla 3}). {[}NO\_CITE{]}

Los resultados confirmaron a la IgG basal elevada y a la eosinofilia
basal como los principales predictores exploratorios de una cinética de
aclaramiento más lenta. {[}NO\_CITE{]}

\subsubsection{Trayectoria paramétrica
(LMM)}\label{trayectoria-paramuxe9trica-lmm}

El modelo lineal mixto (LMM) estimó una pendiente de descenso de
log(IgG) de \textbf{−0,068 unidades/mes} (p \textless{} 0,001),
equivalente a una vida media serológica de \textbf{10,2 meses}
(\textbf{Figura 5}). {[}NO\_CITE{]}

El tiempo estimado de cruce del umbral de negatividad por LMM fue de
\textbf{10,9 meses}, consistente con la estimación no paramétrica de
Kaplan--Meier. {[}NO\_CITE{]}

La interacción tiempo × sexo no fue significativa (p = 0,501),
sugiriendo velocidad de descenso similar entre grupos pese a diferencias
en los títulos basales. {[}NO\_CITE{]}

\subsubsection{Persistencia y recidiva}\label{persistencia-y-recidiva}

Se identificaron \textbf{5 recidivas serológicas} (11,6 \%) y \textbf{12
persistencias} (27,9 \%) sobre 43 pacientes clasificables; no hubo
recidiva parasitológica documentada (\textbf{Tabla 4}). {[}NO\_CITE{]}

Las trayectorias individuales (\textbf{Figura 6}, diagrama de nadador)
mostraron dos patrones: caídas monótonas en curación sostenida frente a
rebotes en ``V'' en recidivas. {[}NO\_CITE{]}

La concordancia entre seronegativización y normalización eosinofílica a
los 12 meses fue nula (κ negativo, \textbf{Tabla 5}), subrayando la
independencia de ambos biomarcadores en el seguimiento post-tratamiento.
{[}NO\_CITE{]}

\subsection*{Discusión}\label{discusiuxf3n}
\addcontentsline{toc}{subsection}{Discusión}

\subsubsection{Principales hallazgos}\label{principales-hallazgos}

Esta cohorte retrospectiva de 44 adultos con estrongiloidiasis tratados
con ivermectina aporta una caracterización cuantitativa de la cinética
de seronegativización por IgG-ELISA en un entorno no endémico.
{[}NO\_CITE{]}

La mediana de seronegativización absoluta (8,8 meses) se sitúa en el
rango inferior-medio de la literatura, que reporta tiempos medianos de 6
a 24 meses según umbral y definición.
{[}CITED:repetto\_2018\_longterm{]}
{[}CITE?:buonfrate\_2018\_followup{]}

La cinética de primer orden estimada por el LMM (pendiente −0,068 u/mes;
vida media 10,2 meses) es consistente con modelos previos de descenso
exponencial de IgG post-antihelmíntico.
{[}CITE?:bisoffi\_2014\_salvador{]}

\subsubsection{IgG basal como predictor (hallazgo
principal)}\label{igg-basal-como-predictor-hallazgo-principal}

El título de IgG basal ≥ mediana se asoció a una reducción del hazard de
seronegativización (HR 0,19; p=0,004). {[}NO\_CITE{]}

Este hallazgo es coherente con estudios previos que sugieren una
asociación inversa entre carga antigénica basal y velocidad de
aclaramiento. {[}CITE?:buonfrate\_2015\_kinetics{]}

La falta de significación tras ajuste BH--FDR refleja el tamaño muestral
limitado (n=39 evaluables) y obliga a interpretar el hallazgo como
generador de hipótesis. {[}NO\_CITE{]}

\subsubsection{Eosinofilia basal como predictor
secundario}\label{eosinofilia-basal-como-predictor-secundario}

La eosinofilia basal emergió como covariable de interés en el modelo
multivariable estratificado por sexo (C = 0,676). {[}NO\_CITE{]}

La literatura describe la eosinofilia como marcador de respuesta Th2
modulable por el huésped, con asociación variable al tiempo de curación
en series pequeñas. {[}CITE?:nutman\_2017\_clinical{]}

La discordancia temporal observada entre normalización eosinofílica
(mediana 3,0 meses) y seronegativización (mediana 8,8 meses) apoya que
ambos biomarcadores capturan fenómenos biológicos distintos.
{[}CITE?:mejia\_2012\_screening{]}

\subsubsection{Persistencia y recidiva}\label{persistencia-y-recidiva-1}

La tasa de persistencia serológica a ≥12 meses (27,9 \%) es comparable a
la descrita en cohortes europeas similares.
{[}CITED:salvador\_2024\_barcelona{]}
{[}CITE?:cutfield\_2024\_auckland{]}

Las 5 recidivas serológicas sin evidencia parasitológica plantean el
dilema clásico entre reinfección silente, fallo terapéutico subumbral y
oscilación técnica del kit ELISA. {[}CITED:repetto\_2018\_longterm{]}

La ausencia de recidiva parasitológica documentada apoya interpretar la
mayoría de rebotes como fluctuación técnica o carga antigénica residual.
{[}NO\_CITE{]}

\subsubsection{Implicaciones clínicas}\label{implicaciones-cluxednicas}

Los datos sugieren que un seguimiento serológico a \textbf{≥ 12 meses}
post-ivermectina es razonable antes de considerar persistencia como
fallo terapéutico. {[}NO\_CITE{]}

Pacientes con IgG basal elevada podrían requerir una ventana de
seguimiento extendida (≥ 18--24 meses) antes de cerrar el caso.
{[}NO\_CITE{]}

La discordancia entre eosinófilos e IgG implica que ambos biomarcadores
deben valorarse de forma independiente, no como marcadores redundantes
de curación. {[}NO\_CITE{]}

\subsubsection{Fortalezas}\label{fortalezas}

\begin{itemize}
\tightlist
\item
  Cohorte monocéntrica con datos longitudinales estandarizados y
  protocolo de seguimiento homogéneo. {[}NO\_CITE{]}
\item
  Análisis pre-especificado con ajuste por comparaciones múltiples,
  evitando cherry-picking. {[}NO\_CITE{]}
\item
  Triangulación no paramétrica (KM) y paramétrica (LMM) de la cinética.
  {[}NO\_CITE{]}
\item
  Evaluación de 7 definiciones operativas del evento (análisis de
  sensibilidad robusto). {[}NO\_CITE{]}
\end{itemize}

\subsubsection{Limitaciones}\label{limitaciones}

\begin{itemize}
\tightlist
\item
  Tamaño muestral pequeño (n=44 basal, 39 evaluables) con baja potencia
  para detectar efectos moderados tras ajuste por comparaciones
  múltiples. {[}NO\_CITE{]}
\item
  Diseño retrospectivo con censura administrativa heterogénea (mediana
  11,7 meses). {[}NO\_CITE{]}
\item
  Cambio de kit ELISA durante el periodo como posible fuente de
  variabilidad técnica. {[}NO\_CITE{]}
\item
  Ausencia de IgG basal en el 18 \% de los pacientes; aunque Little's
  MCAR no fue significativa (p=0,11), no puede descartarse un sesgo
  residual. {[}NO\_CITE{]}
\item
  Cohorte geográficamente sesgada (predominio sudamericano), limitando
  la generalización a pacientes de otras regiones endémicas.
  {[}CITED:salvador\_2020\_systematic{]}
\item
  Imposibilidad de diferenciar reinfección de fallo terapéutico en
  pacientes con exposición persistente o viajes a área endémica.
  {[}NO\_CITE{]}
\end{itemize}

\subsubsection{Comparación con la
literatura}\label{comparaciuxf3n-con-la-literatura}

En la red +REDIVI (n=1245) y en la cohorte de Barcelona (n=865), la
distribución demográfica es comparable (predominio masculino, origen
sudamericano mayoritario). {[}CITED:salvador\_2019\_redivi{]}
{[}CITED:salvador\_2024\_barcelona{]}

Los tiempos de seronegativización reportados en dichas series son
consistentes con nuestros hallazgos, aunque con menor granularidad
estadística. {[}CITE?:gamez\_2022\_spanish{]}

El ensayo clínico Strong Treat (1--4) sobre dosis única vs múltiple de
ivermectina apoya la elección del régimen en nuestra cohorte; la
heterogeneidad basal observada en nuestro estudio justifica la pauta
múltiple en casos seleccionados.
{[}CITED:buonfrate\_2019\_strongtreat{]}

\subsubsection{Propuesta para el protocolo de
seguimiento}\label{propuesta-para-el-protocolo-de-seguimiento}

En base a los hallazgos, proponemos un protocolo simplificado de
seguimiento post-ivermectina en pacientes inmunocompetentes con
estrongiloidiasis confirmada:

\begin{enumerate}
\def\labelenumi{\arabic{enumi}.}
\tightlist
\item
  Control serológico a los 6, 12 y 24 meses post-tratamiento.
  {[}NO\_CITE{]}
\item
  Extender a 18 meses si IgG basal \textgreater{} mediana poblacional.
  {[}NO\_CITE{]}
\item
  Considerar curación definitiva ante dos controles consecutivos con IgG
  \textless{} 1,1. {[}NO\_CITE{]}
\item
  Reevaluar parasitológicamente solo en persistencia clínica o rebote
  serológico sostenido. {[}NO\_CITE{]}
\end{enumerate}

\subsection*{Conclusiones}\label{conclusiones}
\addcontentsline{toc}{subsection}{Conclusiones}

\begin{enumerate}
\def\labelenumi{\arabic{enumi}.}
\item
  La seronegativización por IgG-ELISA tras ivermectina en adultos con
  estrongiloidiasis sigue una cinética de primer orden, con una vida
  media serológica aproximada de 10 meses en un entorno no endémico.
  {[}NO\_CITE{]}
\item
  La mayor parte de los pacientes inmunocompetentes alcanzan la
  negativización serológica dentro de los 24 meses post-tratamiento, lo
  que sitúa la ventana razonable de seguimiento en dicho intervalo.
  {[}NO\_CITE{]}
\item
  La IgG basal elevada y, de forma secundaria, la eosinofilia basal
  emergen como marcadores exploratorios de aclaramiento más lento, con
  implicaciones prácticas para individualizar la duración del
  seguimiento. {[}NO\_CITE{]}
\item
  La discordancia temporal entre normalización eosinofílica y
  seronegativización obliga a valorar ambos biomarcadores de manera
  independiente en el control postratamiento. {[}NO\_CITE{]}
\item
  Los hallazgos, generadores de hipótesis en un piloto monocéntrico
  (n=44), deben replicarse en cohortes prospectivas mayores y
  multicéntricas antes de modificar protocolos clínicos. {[}NO\_CITE{]}
\item
  El trabajo contribuye al cuerpo de evidencia sobre la práctica real
  del seguimiento serológico de estrongiloidiasis importada en España y
  aporta una trazabilidad metodológica reproducible (código y datos
  versionados). {[}NO\_CITE{]}
\end{enumerate}

\newpage

\subsection*{Bibliografía}\label{bibliografuxeda}
\addcontentsline{toc}{subsection}{Bibliografía}

\subsubsection{\texorpdfstring{Lista provisional (generada manualmente
desde
\texttt{preguntas\_opev.md})}{Lista provisional (generada manualmente desde preguntas\_opev.md)}}\label{lista-provisional-generada-manualmente-desde-preguntas_opev.md}

\begin{enumerate}
\def\labelenumi{\arabic{enumi}.}
\tightlist
\item
  \textbf{Salvador F, Treviño B, Chamorro-Tojeiro S, et al.} Imported
  Strongyloidiasis: Data From 1245 Cases Registered in the +REDIVI
  Spanish Collaborative Network (2009--2017). \emph{PLoS Negl Trop Dis.}
  2019;13(5):e0007399. doi:10.1371/journal.pntd.0007399. {[}key:
  \texttt{salvador\_2019\_redivi}{]}
\item
  \textbf{Asundi A, Beliavsky A, Liu XJ, et al.} Prevalence of
  Strongyloidiasis and Schistosomiasis Among Migrants: A Systematic
  Review and Meta-Analysis. \emph{Lancet Glob Health.}
  2019;7(2):e236--e248. doi:10.1016/S2214-109X(18)30490-X. {[}key:
  \texttt{asundi\_2019\_lancetgh}{]}
\item
  \textbf{Salvador F, Treviño B, Bosch-Nicolau P, et al.}
  Strongyloidiasis Screening in Migrants Living in Spain: Systematic
  Review and Meta-Analysis. \emph{Trop Med Int Health.}
  2020;25(3):281--290. doi:10.1111/tmi.13352. {[}key:
  \texttt{salvador\_2020\_systematic}{]}
\item
  \textbf{Salvador F, Treviño B, Sulleiro E, et al.} Epidemiological and
  Clinical Trends of Imported Strongyloidiasis in a Referral
  International Health Unit, Barcelona, Spain. \emph{Travel Med Infect
  Dis.} 2024;58:102690. doi:10.1016/j.tmaid.2024.102690. {[}key:
  \texttt{salvador\_2024\_barcelona}{]}
\item
  \textbf{Cutfield T, Motuhifonua SK, Blakiston M, et al.}
  Strongyloidiasis in Auckland: A Ten-Year Retrospective Study.
  \emph{PLoS Negl Trop Dis.} 2024;18(3):e0012045.
  doi:10.1371/journal.pntd.0012045. {[}key:
  \texttt{cutfield\_2024\_auckland}{]}
\item
  \textbf{Buonfrate D, Bradbury RS, Watts MR, Bisoffi Z.} Human
  Strongyloidiasis: Complexities and Pathways Forward. \emph{Clin
  Microbiol Rev.} 2023;36(4):e0003323. doi:10.1128/cmr.00033-23. {[}key:
  \texttt{buonfrate\_2023\_cmr}{]}
\item
  \textbf{Czeresnia JM, Weiss LM.} \emph{Strongyloides stercoralis.
  Lung.} 2022;200(2):141--148. doi:10.1007/s00408-022-00528-z. {[}key:
  \texttt{czeresnia\_2022\_lung}{]}
\item
  \textbf{Buonfrate D, Salas-Coronas J, Muñoz J, et al.} Multiple-Dose
  Versus Single-Dose Ivermectin for \emph{Strongyloides stercoralis}
  Infection (Strong Treat 1 to 4). \emph{Lancet Infect Dis.}
  2019;19(11):1181--1190. doi:10.1016/S1473-3099(19)30289-0. {[}key:
  \texttt{buonfrate\_2019\_strongtreat}{]}
\item
  \textbf{Repetto SA, Ruybal P, Batalla E, et al.} Strongyloidiasis
  Outside Endemic Areas: Long-Term Parasitological and Clinical
  Follow-Up After Ivermectin Treatment. \emph{Clin Infect Dis.}
  2018;66(10):1558--1565. doi:10.1093/cid/cix1069. {[}key:
  \texttt{repetto\_2018\_longterm}{]}
\item
  \textbf{Benjamini Y, Hochberg Y.} Controlling the False Discovery
  Rate: a Practical and Powerful Approach to Multiple Testing. \emph{J R
  Stat Soc Ser B.} 1995;57(1):289--300. {[}key:
  \texttt{benjamini\_1995\_fdr}{]}
\item
  \textbf{von Elm E, Altman DG, Egger M, et al.} The Strengthening the
  Reporting of Observational Studies in Epidemiology (STROBE) Statement.
  \emph{Lancet.} 2007;370(9596):1453--1457.
  doi:10.1016/S0140-6736(07)61602-X. {[}key:
  \texttt{vonelm\_2007\_strobe}{]}
\item
  \textbf{Inagaki K, Bradbury RS, Hobbs CV.} Hospitalizations Associated
  With Strongyloidiasis in the United States, 2003--2018. \emph{Clin
  Infect Dis.} 2022;75(9):1548--1555. doi:10.1093/cid/ciac220. {[}key:
  \texttt{inagaki\_2022\_hospitalizations}{]}
\end{enumerate}

\newpage

\subsection*{Agradecimientos}\label{agradecimientos}
\addcontentsline{toc}{subsection}{Agradecimientos}

A los equipos de Medicina Interna, Enfermedades Infecciosas y
Microbiología del Hospital Central de la Defensa ``Gómez Ulla'' por el
mantenimiento de las bases de datos clínicas y serológicas que han hecho
posible este trabajo. {[}NO\_CITE{]}

Al Centro Nacional de Microbiología (Instituto de Salud Carlos III) por
la realización de la serología IgG-ELISA durante el periodo de estudio.
{[}NO\_CITE{]}

A {[}TUTOR/A PENDIENTE{]} por la dirección académica del trabajo y la
revisión crítica del manuscrito. {[}NO\_CITE{]}

A la Facultad de Medicina y Ciencias de la Salud de la Universidad de
Alcalá por el marco formativo del Trabajo Fin de Grado. {[}NO\_CITE{]}

Los autores declaran no tener conflictos de interés. El estudio no
recibió financiación externa. {[}NO\_CITE{]}

\newpage

\subsection*{Anexo académico --- Reflexión para la evaluación del
TFG}\label{anexo-acaduxe9mico-reflexiuxf3n-para-la-evaluaciuxf3n-del-tfg}
\addcontentsline{toc}{subsection}{Anexo académico --- Reflexión para la
evaluación del TFG}

\emph{Esta sección pertenece exclusivamente a la versión académica del
trabajo (TFG) y no se incluirá en la versión enviada a revistas
científicas.}

\subsubsection{Competencias movilizadas}\label{competencias-movilizadas}

\begin{itemize}
\tightlist
\item
  \textbf{CB1--CB5} (competencias básicas del título): integración de
  conocimientos clínicos (helmintiasis importada), metodología
  epidemiológica (STROBE, Kaplan--Meier, Cox, LMM) y difusión científica
  (manuscrito estructurado según normas ICMJE). {[}NO\_CITE{]}
\item
  \textbf{CG1, CG7} (generales): capacidad de análisis crítico de la
  literatura y aplicación de modelos estadísticos a datos reales.
  {[}NO\_CITE{]}
\item
  \textbf{CE1, CE5, CE10} (específicas de Medicina): razonamiento
  diagnóstico, manejo farmacológico (ivermectina) y diseño de protocolos
  de seguimiento. {[}NO\_CITE{]}
\end{itemize}

\subsubsection{Reflexión
metodológica}\label{reflexiuxf3n-metodoluxf3gica}

El trabajo ha requerido equilibrar la exigencia estadística
(comparaciones múltiples, modelos mixtos) con las limitaciones de tamaño
muestral de una cohorte monocéntrica (n=44). {[}NO\_CITE{]}

Se ha adoptado un enfoque de \textbf{estudio piloto exploratorio},
evitando el cherry-picking metodológico y manteniendo la distinción
entre endpoint primario pre-especificado (seronegativización absoluta) y
análisis de sensibilidad (umbrales alternativos, confirmación repetida,
biomarcadores secundarios). {[}NO\_CITE{]}

La reproducibilidad se ha priorizado mediante versionado del código
(Python + Quarto) y trazabilidad de decisiones (\texttt{prompt.md},
auditorías). {[}NO\_CITE{]}

\subsubsection{Limitaciones como
aprendizaje}\label{limitaciones-como-aprendizaje}

\begin{itemize}
\tightlist
\item
  La gestión de datos ausentes (18 \% IgG basal) ha obligado a
  considerar MCAR/MAR y la eventual aplicación de imputación múltiple
  (MICE), descartada por impacto marginal sobre el endpoint primario.
  {[}NO\_CITE{]}
\item
  La heterogeneidad técnica (cambio de kit ELISA) ha forzado un análisis
  de sensibilidad adicional no anticipado. {[}NO\_CITE{]}
\item
  La escasa potencia estadística tras ajuste BH--FDR es una lección
  sobre la planificación a priori del tamaño muestral en estudios de
  carácter exploratorio. {[}NO\_CITE{]}
\end{itemize}

\subsubsection{Líneas de continuidad}\label{luxedneas-de-continuidad}

\begin{itemize}
\tightlist
\item
  Ampliación de la cohorte a multicéntrica con incorporación de
  inmunodeprimidos (no infrarrepresentados) para valorar el efecto de la
  inmunosupresión sobre la cinética de aclaramiento. {[}NO\_CITE{]}
\item
  Validación prospectiva de la IgG basal como predictor y exploración de
  covariables no evaluadas (carga parasitaria cuantitativa por qPCR,
  subclases de IgG). {[}NO\_CITE{]}
\item
  Evaluación de la concordancia entre kits ELISA (Strongyloides ratti
  antigen-based vs antígenos recombinantes). {[}NO\_CITE{]}
\end{itemize}

\subsubsection{Autoevaluación}\label{autoevaluaciuxf3n}

El grado de consecución de los objetivos pre-especificados ha sido alto
para los objetivos descriptivos (KM, LMM) y moderado para los
predictivos (Cox con ajuste BH--FDR), con resultados consistentes pero
limitados por tamaño muestral. {[}NO\_CITE{]}

La principal aportación personal es la integración de un pipeline
reproducible (datos + código + manuscrito bajo control de versiones) que
soporta futuras ampliaciones y auditorías. {[}NO\_CITE{]}




\end{document}
